\chapter{PCSEL 是什么，为什么重要}
\label{ch:introduction}

\section*{学习目标}
\begin{itemize}
  \item 从激光器应用需求出发理解 PCSEL 为何会成为一个独立方向。
  \item 准确区分 PCSEL、边发射激光器、VCSEL、DFB 激光器与一般光子晶体激光器。
  \item 建立本书的基本视角：PCSEL 同时是开放周期电磁系统、半导体有源腔与真实器件。
\end{itemize}

\section*{先修知识}
基础电动力学、傅里叶光学、半导体激光器基本概念。

\section*{本章路线图}
本章先从“为什么需要这样一种激光器”出发，再说明 PCSEL 的物理图像，随后通过与相邻器件对比划清边界，最后给出贯穿全书的三重视角与关键观测量。

\section{先问问题：为什么会需要 PCSEL}
如果读者已经接触过边发射激光器（edge-emitting laser）或垂直腔面发射激光器（vertical-cavity surface-emitting laser, VCSEL），很容易产生一个疑问：现有半导体激光器种类已经很多，为什么还需要专门讨论 PCSEL？

这个问题不能只从“有没有新结构”来回答，而应从光源需求来回答。许多应用同时希望得到以下几件事：
\begin{itemize}
  \item 面发射，便于阵列集成、封装与自由空间耦合；
  \item 大出光面积，以承受更高功率并降低局部热负荷；
  \item 接近单横模或至少主模占优，以获得高亮度（brightness，单位发射面积、单位立体角内的辐射功率，单位 W·cm$^{-2}$·sr$^{-1}$；提高亮度要求在提升功率的同时收紧远场发散角）和良好光束质量；
  \item 可工程化的远场和偏振，以适配照明、探测、投影或相干系统。
\end{itemize}

这四个目标并不天然相容。出光面积做大以后，横向可容纳的模式数迅速增加；模式一多，远场会恶化，亮度会下降，偏振也更难稳定。边发射激光器通常在一维波导方向上建立反馈，VCSEL 则主要依赖垂直 DBR 腔。它们各有优势，但面对“大面积、面发射、同时保持较好模式控制”这一组目标时，都会碰到各自的边界。PCSEL 正是在这个背景下出现的：它试图利用二维光子晶体的面内分布式反馈，一边把横向模式组织起来，一边允许能量向法向辐射输出。

\begin{IntuitionBox}
PCSEL 的问题意识不是“如何在光子晶体里做出激光”，而是“如何让一个大面积半导体有源区，在面发射条件下仍由少数目标模稳定主导输出”。这句话比任何一句定义都更接近 PCSEL 的工程本质。
\end{IntuitionBox}

\section{PCSEL 的基本图像：面内反馈与垂直出射并存}
\cref{fig:pcsel-layer} 给出了 PCSEL 最基本的层状图像。外延层沿 $z$ 方向堆叠，光子晶体周期主要位于平面内 $(x,y)$。从纵向看，它像一个有源 slab 波导；从横向看，它又像一个二维周期散射网络。正是这两种视角的叠加，构成了 PCSEL 的特殊性。

\begin{figure}[htbp]
  \centering
  \begin{tikzpicture}[x=1cm,y=1cm]
    \draw[fill=gray!15] (0,0) rectangle (8,0.6);
    \draw[fill=blue!15] (0,0.6) rectangle (8,1.2);
    \draw[fill=green!15] (0,1.2) rectangle (8,1.8);
    \draw[fill=yellow!20] (0,1.8) rectangle (8,2.6);
    \draw[fill=red!15] (0,2.6) rectangle (8,3.2);
    \draw[dashed] (0,2.2)--(8,2.2);
    \foreach \x in {0.6,1.4,...,7.4}{
      \draw[fill=white] (\x,2.2) circle (0.13);
    }
    \draw[->,thick] (4,3.2)--(4,4.2) node[above]{法向输出};
    \draw[->,thick,blue!70!black] (2.2,2.2)--(3.4,2.2);
    \draw[->,thick,blue!70!black] (5.8,2.2)--(4.6,2.2);
    \node[left] at (0,0.3) {衬底};
    \node[left] at (0,0.9) {下包层};
    \node[left] at (0,1.5) {波导/有源层};
    \node[left] at (0,2.2) {光子晶体层};
    \node[left] at (0,2.9) {上包层/接触相关层};
  \end{tikzpicture}
  \caption{PCSEL 的纵向层状结构与面内反馈、法向辐射的概念图。}
  \label{fig:pcsel-layer}
\end{figure}

对初学者而言，最容易把 PCSEL 想成“一个有孔的出光面”。这种想法太静态，也太局部。更准确的图像是：有源区中的导模在平面内传播，遇到二维周期折射率调制后发生布拉格散射（Bragg scattering），多个传播方向的波被相互耦合；在某些特定对称条件下，部分面内模式又能与自由空间辐射通道耦合，于是形成法向出光。换句话说，PCSEL 并不是先有一个封闭腔、再在顶部开孔把光漏出去；它从一开始就是一个“面内耦合与向外辐射同时被设计”的腔。

\section{定义 PCSEL：什么必须成立，什么不应混淆}
为了避免和相近器件混淆，本书采用如下工作定义。

\begin{definition}
光子晶体面发射激光器（photonic-crystal surface-emitting laser, PCSEL）是指满足以下特征的一类半导体激光器：
\begin{enumerate}
  \item 主要反馈机制来自平面内二维周期结构引起的分布式耦合，而非单纯的垂直 DBR 腔或单轴波导端面反射；
  \item 主要输出方向为法向或近法向面发射；
  \item 激射模式的存在与选择与布洛赫模（Bloch mode）、$\Gamma$ 点附近模、对称性和辐射耦合直接相关；
  \item 器件性能不能只由被动冷腔模式决定，而必须与外延层、有源区、电注入和热效应联立理解。
\end{enumerate}
\end{definition}

这个定义里最重要的是最后一条。很多文章在介绍 PCSEL 时，会把重点放在光子晶体和能带上。这当然重要，但如果只看到“二维周期共振”，就会把 PCSEL 误看成一类特殊的光子晶体激光器，而忽略它作为半导体器件的本质。对于真正可工作的电注入 PCSEL，纵向层栈、电流路径、载流子输运、热源分布和掺杂损耗，都不是“后处理修正项”，而是决定器件成败的主因素。

\begin{RigorousBox}
本书会反复坚持三个区分：被动共振不等于有源阈值，阈值增益不等于阈值电流，无限周期结果不等于有限器件工作结果。只要把这三个区分守住，很多概念性错误就能避免。
\end{RigorousBox}

\section{与相邻器件逐一比较：边界要靠对比才能真正清楚}
如果只背定义，读者仍然容易在具体结构前犹豫“这到底算不算 PCSEL”。因此，更有效的办法是把 PCSEL 与几类常见器件逐一比较。\cref{tab:device-compare} 总结了最核心的区别。

\begin{table}[htbp]
  \centering
  \caption{PCSEL 与几类相关半导体激光器的核心比较。}
  \label{tab:device-compare}
  \small
  \setlength{\tabcolsep}{4pt}
  \begin{tabularx}{\textwidth}{@{}p{0.17\textwidth}p{0.22\textwidth}p{0.17\textwidth}Y@{}}
    \toprule
    器件类型 & 主要反馈机制 & 主要出光方向 & 模选择的主控因素 \\
    \midrule
    边发射激光器 & 波导端面/DFB/DBR 轴向反馈 & 沿波导方向 & 波导横模、纵模与端面/光栅条件 \\
    VCSEL & 垂直 DBR 腔 & 法向面发射 & 垂直腔谐振、横向氧化孔或电流孔径 \\
    DFB 激光器 & 一维光栅分布式反馈 & 多为边发射 & 一维布拉格条件与相移设计 \\
    一般光子晶体激光器 & 缺陷局域或局域共振 & 可多样 & 局域缺陷模或纳米腔共振 \\
    PCSEL & 二维周期结构面内反馈并与辐射耦合 & 法向或近法向面发射 & $\Gamma$ 点附近模式、二维对称性、辐射耦合、外延与注入分布 \\
    \bottomrule
  \end{tabularx}
\end{table}

这里尤其要强调两点。

在强调之前先补一个常被忽略的器件细节：VCSEL 中的氧化孔（oxide aperture）通过选择性氧化高 Al 组分 AlGaAs 形成绝缘环，其开口直径通常为 \SIrange{3}{20}{\micro\meter}，把有效增益区压缩在小孔径内，从而抑制高阶横模，本质上是空间选择性增益。PCSEL 没有这样的物理孔径，它的横向模式控制依赖二维布拉格散射在大面积孔阵上的相干耦合，模式由散射网络的对称性与损耗分布决定，而不是几何孔径决定。

第一，PCSEL 不是 VCSEL 的“打孔版”。VCSEL 的核心是垂直腔，PCSEL 的核心是平面内耦合。两者都可能面发射，但形成面发射的物理路径不同。VCSEL 中，模式主要先在垂直方向被 DBR 限定；PCSEL 中，模式先要在平面内通过周期散射建立特定的反馈结构，再通过允许的辐射通道向外输出。

第二，PCSEL 也不是一般“photonic crystal laser”的同义词。很多光子晶体激光器依赖的是小体积极强局域模式，而 PCSEL 的一个重要目标恰恰是做大有效出光面积，同时尽可能维持良好的模式控制。一个靠纳米缺陷局域实现激射的光子晶体纳米腔，原则上不应自动归入 PCSEL。

\begin{ExampleBox}
如果一个结构具有二维周期孔阵，但主要反馈实际上来自上下 DBR 形成的垂直腔，而孔阵只是辅助出光或横向整形，那么把它称为 PCSEL 就需要非常谨慎。是否属于 PCSEL，不看表面形状，而看主要反馈是谁提供的。
\end{ExampleBox}

\section{为什么大面积单模困难，而 PCSEL 可能有优势}
对于任何激光器，只要横向尺寸变大，允许存在的横向模式数目就会增加。若没有额外机制筛选模式，不同模式的阈值会变得接近，结果就是：
\begin{itemize}
  \item 起振模式不稳定，容易随注入变化切换；
  \item 副瓣增强，远场变差；
  \item 虽然总功率升高，但亮度（brightness）和光束质量（beam quality）不一定改善。
\end{itemize}

这里的亮度不是总功率的同义词，而是前文已经定义过的“单位发射面积、单位立体角内的辐射功率”。也就是说，器件即便输出总功率更高，只要发光面积显著变大或远场发散角明显变宽，亮度仍可能下降。对 PCSEL 来说，真正困难的目标是同时扩大有效出光面积、压住远场发散角，并尽量维持单模与单偏振工作。

定量上常用近似写法是
\begin{equation}
B \approx \frac{P}{A_{\mathrm{eff}}\,\Omega},\qquad \Omega\approx \theta_x\theta_y,
\end{equation}
其中 $P$ 为输出功率，$A_{\mathrm{eff}}$ 为有效出光面积，$\Omega$ 为远场立体角（小角度时可用发散半角 $\theta_x,\theta_y$ 近似）。这一定义提醒我们：亮度提升必须同时压缩发散角或降低等效束腰面积，不能只靠总功率上升。

PCSEL 的潜在优势在于，它并不是先做大面积，再事后用损耗把高阶模压下去；相反，它可以在结构层面把模式选择机制提前嵌入平面内反馈网络中。二维光子晶体带来的不是“任意复杂”，而是“可工程化的耦合通道”：通过晶格常数、孔形、扰动、双晶格、对称性破缺等设计，研究者可以调节模式分裂、辐射方向性和偏振选择规则。这正是 PCSEL 能在大面积单模方向上持续吸引研究的原因。

但这里必须把“可能有优势”这句话说得更硬一些。PCSEL 的大面积单模潜力来自\textbf{模式组织机制}，而不是天然保证。更准确地说，PCSEL 的大面积单模能力不是 band diagram 单独的性质，而是
\[
\text{band}
\;+\;
\text{radiation}
\;+\;
\text{finite size}
\;+\;
\text{injection}
\;+\;
\text{thermal}
\]
这五类因素共同决定的联合性质。只要其中任意一环失控，例如有限阵列边缘引入新泄漏通道、电流拥挤把增益峰拖离目标模，或者温升把模式身份重新排序，那么“名义上看起来很优的候选模”就未必能在器件里兑现成稳定单模工作。

但这里有一个必须提前说清楚的事实：PCSEL 只是“有潜力”，不是“天然单模”。如果外延层设计不当、电流拥挤严重、热分布失控，那么再漂亮的能带和 $\Gamma$ 点模式也可能在器件层面失效。也就是说，PCSEL 的优势只有在光学、电学和热学同时被兼顾时才会兑现。

\begin{PitfallBox}
初学者常把“看到一个高 $Q$ 的 $\Gamma$ 点候选模”理解成“这就是低阈值大面积单模器件”。这一步跨得太大了。高 $Q$ 只说明被动泄露小；器件是否低阈值、是否稳定单模，还要看有源区重叠、内部损耗、注入非均匀和热反馈。
\end{PitfallBox}

\section{PCSEL 必须同时用三种视角来理解}
从本章开始，本书会一直交替使用三种视角。这三种视角不是三套互不相干的话语，而是一条连续推理链上的三个层面。

\subsection{第一种视角：开放周期电磁系统}
从这一视角看，PCSEL 是一个与自由空间耦合的二维开放周期系统。关键词是 Maxwell 方程、Bloch 理论、$\Gamma$ 点、辐射损耗、复频率、远场和偏振选择。这个层面回答的问题包括：有哪些候选模？哪些模能向法向辐射？辐射损耗大小规律是什么？

\subsection{第二种视角：半导体有源光学腔}
从这一视角看，PCSEL 不只是“有模”，而是“有需要被增益补偿的损耗”。关键词是材料增益、模增益、光学约束因子、内部损耗、阈值增益和模式竞争。这个层面回答的问题包括：有源区与目标模重叠如何？哪一个模式更容易先达到光学阈值？

\subsection{第三种视角：真实电注入器件}
从这一视角看，PCSEL 不再只是光场问题，而是一个由外延层、电流扩展、载流子输运、Joule 热和热致折射率变化共同决定的器件。关键词是漂移-扩散（drift-diffusion）、Poisson 方程、热方程、注入效率、电流拥挤、温升和热-电-光耦合。这个层面回答的问题包括：阈值电流是多少？温升会不会改变模式排序？器件在多大电流窗口内能稳定工作？

这三种视角之间的关系可以概括为：第一层决定候选模，第二层决定光学阈值排序，第三层决定器件真正能否达到并维持该排序。更具体的接口链条可以用一条最小的量级传递线来记：
\begin{equation}
Q,\tilde{\omega}
\;\Rightarrow\;
\alpha_{\mathrm{rad}}\,(=\omega/2Qv_g)
\;\Rightarrow\;
g_{\mathrm{th}}=\alpha_{\mathrm{rad}}+\alpha_{\mathrm{int}}
\;\Rightarrow\;
I_{\mathrm{th}}\ (\text{由 }g_{\mathrm{th}}\text{ 进入载流子/热方程}).
\end{equation}
它提醒读者：$Q$ 与辐射损耗是 L1 的输出，阈值增益是 L2 的入口，阈值电流与工作窗口则必须进入 L3/L4 的输运与热回写。

\section{关键观测量：本书后面所有章节都在回答哪些量}
本书后文经常出现的观测量并不是相互独立的一串名词，它们共同刻画了 PCSEL 的工作状态。

\begin{itemize}
  \item 品质因子 $Q$ 与光子寿命：属于被动腔信息，用来表征损耗时间尺度；
  \item 阈值增益 $g_\mathrm{th}$：属于有源光学层，回答“补偿损耗至少需要多少模增益”；
  \item 阈值电流 $I_\mathrm{th}$：属于器件层，回答“真实器件要打多大电流才能达到阈值”；
  \item 远场与偏振：回答输出是否在目标方向、目标偏振上集中；
  \item 线宽、相干性、亮度与 $M^2$：回答器件作为系统光源是否真正可用。
\end{itemize}

把这些量按层级区分非常重要。比如，$Q$ 很高不自动意味着 $I_\mathrm{th}$ 很低；远场很好看也不自动意味着器件热稳定；有源区放进去了也不自动意味着目标模能拿到足够模增益。后文每一章都会回到这些量，但会明确说明自己讨论的是哪一个层级。

\begin{table}[htbp]
  \centering
  \footnotesize
  \caption{本书最常用观测量的“物理量 $\rightarrow$ 模型层级 $\rightarrow$ 求解器 $\rightarrow$ 实验量”总映射。表中“求解器”表示典型接口，而非唯一实现。}
  \label{tab:quantity-model-solver-exp}
  \setlength{\tabcolsep}{4pt}
  \begin{tabularx}{\textwidth}{@{}p{2.2cm}p{2.5cm}p{3.3cm}Y@{}}
    \toprule
    物理量 & 主要模型层级 & 常见求解器或方法接口 & 典型实验对应量 \\
    \midrule
    带图、Gamma 点候选模 & L1 被动电磁层 & PWEM、CWT、RCWA 的单胞筛选 & 反射/透射谱中的候选特征、偏振选择趋势 \\
    被动 $Q$、辐射损耗、上下出光比例 & L1 被动电磁层 & RCWA、FDTD、FEM、周期频域本征/散射模型 & 线宽收缩趋势、反射谷宽度、远场与上下出光比 \\
    阈值增益 $g_{\mathrm{th}}$ & L2 有源光学层 & CWT、RCWA 阈值代理、重叠积分模型 & 起振前后增益需求、阈值附近谱特征变化 \\
    阈值电流 $I_{\mathrm{th}}$ & L3/L4 器件层 & Poisson + 漂移扩散 + 热回写 & L--I 曲线拐点、J--V--L、热滚降起点 \\
    远场、偏振、$M^2$ & L1/L4 之间的输出层 & RCWA、FDTD、FEM、近远场变换 & CCD 远场、偏振仪、束质测量 \\
    线宽、相干性、频率噪声 & 动态/噪声层 + 实验闭环 & 速率方程、降阶噪声模型、热灵敏度模型 & 频谱仪、自外差/拍频、频率噪声谱测量 \\
    载流子分布、温度场 & L3/L4 器件层 & 半导体电学求解器、热传导求解器（如 CHARGE、Semiconductor、HEAT、Heat Transfer） & 热阻、热像、I--V、谱漂移、阈值随温度变化 \\
    \bottomrule
  \end{tabularx}
\end{table}

\section{本书的主线：从 Maxwell 到工作器件}
如果只用一句话概括全书的结构，可以写成下面这条链：
\begin{equation}
\begin{aligned}
\text{开放 Maxwell 方程}
&\rightarrow
\text{Bloch 模与 }\Gamma\text{ 点}
\rightarrow
\text{耦合与辐射}
\\
&\rightarrow
\text{阈值增益}
\rightarrow
\text{电注入与热回写}
\rightarrow
\text{器件工作窗口}.
\end{aligned}
\end{equation}

这个顺序不是写作上的方便，而是逻辑上的必要。要谈模式，先要知道电磁问题是什么；要谈阈值，先要知道损耗从哪里来；要谈阈值电流，先要知道载流子如何进入有源区；要谈最终远场和稳定性，先要知道温升如何回写光学性质。读者如果沿着这条链去读，就不会把局部结果误当全局结论。

\section*{本章小结}
PCSEL 值得单独成书，不是因为它多了一个“光子晶体”名词，而是因为它把二维周期电磁学、半导体有源腔理论和真实器件电热问题压到了同一个对象中。真正理解 PCSEL，必须从一开始就接受这三重身份。

从下一章起，本书先暂停物理叙述，用一章时间统一坐标系、时间因子、模型层级和常用符号。这一步看似管理性质的，却是全书推导语义不出错的前提——读者在这里对约定达成一次共识，后面就不必在每个公式前先猜“虚部是增益还是损耗”。

\section*{练习题}
\begin{enumerate}
  \item \basicq 不查资料，用自己的语言说明：为什么“大面积面发射”与“单模高亮度”之间天然存在张力？

  \item \basicq 试着分别用一句话概括边发射激光器、VCSEL 和 PCSEL 的主要反馈机制。

  \item \midq 给出一个你认为“不应称为 PCSEL”的结构例子，并说明理由。

  \item \hardq 解释为什么“高 $Q$ 候选模”不能直接推出“低阈值电流器件”。
\end{enumerate}

\section*{建议继续阅读}
\begin{itemize}
  \item \cite{joannopoulos2008} 中关于周期介质与光子晶体基本概念的章节。 这条阅读的重点是补足本章关于 PCSEL 的三重身份与关键观测量的标准背景、经典语言和代表性文献接口。

  \item \cite{coldren2012} 中关于半导体激光器工作机制与器件指标的章节。 这条阅读的重点是补足本章关于 PCSEL 的三重身份与关键观测量的标准背景、经典语言和代表性文献接口。
\end{itemize}


